\documentclass[12pt]{article}
\usepackage{amsmath}
\usepackage{amssymb}
\usepackage{graphicx}
\usepackage{hyperref}
\usepackage{geometry}
\usepackage{color}

\geometry{a4paper, margin=1in}

\title{Laboratory: The Mind Reader (Markov Chains \& Statistics)}
\author{Escher Droste Laboratory Series}
\date{}

\begin{document}
\maketitle

\section*{Introduction}
Humans are remarkably poor at generating true randomness. When asked to simulate a coin flip by typing `1' or `2', the human brain attempts to artificially ``balance'' the sequence, actively avoiding long streaks and falling into predictable, alternating patterns.

In this laboratory, you will build a mathematical ``Mind Reader'' that exploits this human predictability. By implementing Markov Chains of increasing complexity, you will build an Artificial Intelligence capable of guessing a human's next keystroke with extraordinary accuracy. You will then use statistical hypothesis testing to prove the AI has ``learned'' something non-random, before finally destroying its predictive power using the pure mathematical randomness of $\pi$.

All of your work will be done in \texttt{implementation\_tasks.py}. You can add your own \lstinline|unittest.TestCase| classes at the bottom and run \lstinline|python implementation_tasks.py| to execute them.

\subsubsection*{Self-Testing (Optional)}
At the bottom of \lstinline|implementation_tasks.py|, you will find a block that looks like this:
\begin{lstlisting}[language=Python]
if __name__ == '__main__':
    import unittest
    # Add your own unittest.TestCase classes here...
    unittest.main(verbosity=2)
\end{lstlisting}
This allows you to write your own automated unit tests to verify your code before running the visual dashboard. If you'd like to test your functions with specific inputs, simply define a class inheriting from \lstinline|unittest.TestCase| above the \lstinline|unittest.main()| call. For example, if you wanted to test simple math:
\begin{lstlisting}[language=Python]
class TestMyMath(unittest.TestCase):
    def test_addition(self):
        result = 2 + 2
        self.assertEqual(result, 4, "Addition is broken!")
\end{lstlisting}
You can run these tests anytime by executing \lstinline|python implementation_tasks.py| in your terminal. This is completely optional, but highly recommended for debugging tricky edge cases.
\section{Level 1: The Basic N-Gram (Fixed Memory)}
A Markov Chain is a stochastic model describing a sequence of possible events where the probability of each event depends only on the state attained in the previous events.

If we use a fixed memory length $N$ (an $N$-Gram model), the ``state'' is simply the exact sequence of the last $N$ keystrokes.
Your first task is to implement the core mechanics of a fixed-order Markov Chain.

\textbf{Tasks:}
\begin{enumerate}
    \item Implement \texttt{update\_markov\_counts(history, new\_char, order, counts\_dict)} to extract the preceding state string of length $N$ and update the frequency dictionary.
    \item Implement \texttt{predict\_fixed\_order(history, order, counts\_dict)}. Look up the current state in the dictionary and return the character (`1' or `2') with the higher historical frequency. If the state is unseen or perfectly tied, guess randomly.
\end{enumerate}

\section{Level 2: Statistical Significance (Hypothesis Testing)}
When you run your Level 1 predictor, you might see an accuracy of 60\%. But is that mathematically significant, or just a lucky streak?

Under the Null Hypothesis (assuming the human is perfectly random), your sequence of correct guesses $k$ out of $n$ trials follows a Binomial distribution $B(n, 0.5)$. For large $n$, this approximates to a Gaussian (Normal) distribution:
$$ \mathcal{N}\left(\mu = 0.5n, \sigma = 0.5\sqrt{n}\right) $$

\textbf{Task:}
\begin{enumerate}
    \item Implement \texttt{calculate\_p\_value(correct\_guesses, total\_trials)}. Use \texttt{scipy.stats.norm.sf} (the survival function) to calculate the exact p-value---the probability that a perfectly random sequence would achieve your score or higher.
\end{enumerate}
\textit{Watch the red line on the UI. If your p-value drops below 0.05, you have statistically proven the human is predictable!}

\section{Level 3: Variable-Order Fallback (Smoothing)}
The weakness of a fixed $N$-Gram is data sparsity. If $N=5$, the model must see the exact same 5-character sequence before it can make an informed prediction. If it hasn't, it blindly guesses.

We can fix this using a technique similar to \textit{Prediction by Partial Matching}. 

\textbf{Task:}
\begin{enumerate}
    \item Implement \texttt{predict\_with\_fallback(history, max\_order, counts\_dict)}. Instead of rigidly looking at length $N$, loop backwards from $N$ down to 1. If a state has zero counts, or its predictions are perfectly tied, \textit{fall back} to the shorter history length.
\end{enumerate}

\section{Level 4: Mixture of Experts (Online Learning)}
Even with fallback, how do we know the optimal maximum memory length? A short memory adapts quickly but misses complex patterns (High Bias). A long memory catches complex patterns but is easily fooled by noise (High Variance).

Instead of choosing, we will run $M$ different Markov models in parallel (Orders 1 through $M$) as competing ``Experts''. We will track their recent accuracy using an Exponential Decay factor $\gamma$.
$$ \text{Score}_{i} = \text{Score}_{i} \times \gamma + \begin{cases} 1 & \text{if Expert } i \text{ guessed correctly} \\ 0 & \text{otherwise} \end{cases} $$

\textbf{Tasks:}
\begin{enumerate}
    \item Implement \texttt{update\_expert\_scores(scores, expert\_predictions, actual\_char, decay)}.
    \item Implement \texttt{get\_best\_expert(scores)} using \texttt{numpy.argmax} to dynamically route the final prediction to the currently winning Expert.
\end{enumerate}
\textit{Watch the Expert Dashboard on the UI. The AI will dynamically shift gears, using Order 1 for simple sequences and Order 5 for tricky ones!}

\section{Level 5: The $\pi$ Simulator (True Randomness)}
We have proven humans are predictable. But how does our AI perform against strict mathematical pseudo-randomness?
We will generate binary digits of $\pi$ using Machin's arctangent formula:
$$ \frac{\pi}{4} = 4 \arctan\left(\frac{1}{5}\right) - \arctan\left(\frac{1}{239}\right) $$

\textbf{Task:}
\begin{enumerate}
    \item Implement \texttt{generate\_pi\_binary\_digits(num\_bits)}. Use the provided \texttt{arccot} integer-arithmetic helper to calculate $num\_bits$ of $\pi$. Strip the integer part and return the binary string.
\end{enumerate}
\textit{Press the ``Test Against $\pi$'' button. Watch the Markov Chain's accuracy violently collapse to exactly 50\%, trapping the red line in the dead center of the Bell Curve!}

\end{document}
